% !PRINTED 218
\weberpage{218}
Among the roots of the division equation there occur also the
$\varphi(n)$ quantities

\weq{$R_1$}{\mathrm{sn}\,h\,\Omega_1}
which, when $h$ runs through a complete system of numbers incongruent
and coprime to $n$, are all different from one another. The system
$(R_1)$ we shall call the \gesperrt{first row of the roots}.

If now $\mathrm{sn}\,\Omega_2$ is a root not contained in $(R_1)$, then the
$\varphi(n)$ quantities

\weq{$R_2$}{\mathrm{sn}\,h\,\Omega_2,}
which are different both among themselves and from the roots $(R_1)$,
form a second row; and in this way one can proceed until all the
$\varphi(n)\psi(n)$ roots of the proper division equation are distributed
into $\psi(n)$ rows of $\varphi(n)$ terms each.

By the multiplication theorem, every other root of the same row can be
expressed \gesperrt{rationally} by means of one root, in the form:

\weq{1}{\mathrm{sn}\,h\,\Omega = f_h(\mathrm{sn}\,\Omega),}
where $f_h$ is a rational function depending only on the multiplier $h$,
not on the choice of $\Omega$.

If the two roots $x_{\mu,\mu'}$, $x_{\nu,\nu'}$ belong to the same row, then the number $h$ must admit of being so determined that

\weq{2}{h\mu \equiv \nu, \quad h\mu' \equiv \nu' \pmod{n},}
and conversely, if this is the case, then the two roots belong to the
same row. From (2), however, there follows the congruence

\weq{3}{\mu\nu' - \nu\mu' \equiv 0 \pmod{n},}
and from this, conversely, the congruences (2) can be derived again.

For since $\mu, \mu', n$ are relatively prime, one can determine the
numbers $\alpha, \beta$ so that
\[
\alpha\mu - \beta\mu' \equiv 1 \pmod{n}
\]
holds, and from this there follows by means of (3):
\[
\nu \equiv (\alpha\nu - \beta\nu')\mu, \quad \nu' \equiv (\alpha\nu - \beta\nu')\mu' \pmod{n},
\]
hence, if $\alpha\nu - \beta\nu' = h$ is put, the congruences (2).

\gesperrt{The congruence (3) is therefore the necessary and sufficient
condition that the two roots of the proper
division equation $x_{\mu,\mu'}$, $x_{\nu,\nu'}$ belong to the same
row.}
